% !TeX root = ../main.tex

Ein Teil der gewonnenen Erkenntnisse lässt sich auf andere Anwendungsszenarien übertragen.
Die grundsätzliche Erkenntnis, dass Retry-Count und Delay einen Kompromiss zwischen Erfolgsrate und Ausführungszeit
erzeugen, ist nicht auf eine einzelne Systemkonfiguration beschränkt.
Auch die Beobachtung, dass die Wahl der Retry-Strategie von der jeweiligen Fehlerklasse abhängt, stellt eine
grundsätzliche Erkenntnis der Untersuchung dar.
Die konkreten Ergebnisse für Deadlocks, Serialization Failures und Lock Timeouts sind hingegen an das verwendete DBMS
und dessen Verhalten gebunden.
Andere DBMS können Konflikte unterschiedlich behandeln und dadurch andere Ergebnisse und optimale Retry-Konfigurationen
aufweisen.

Für weiterführende Untersuchungen bietet sich zunächst eine Betrachtung unter realer beziehungsweise realitätsnaher
Last an.
Dadurch könnte untersucht werden, ob die beobachteten Zusammenhänge zwischen Erfolgsrate, Ausführungszeit und
Retry-Konfiguration auch unter produktionsähnlichen Bedingungen bestehen.
Darüber hinaus könnte der untersuchte Parameterbereich erweitert werden.
Weitere Delay und Retry-Werte könnten zeigen, ob die in dieser Untersuchung ermittelten Konfigurationen tatsächlich
günstige Bereiche darstellen oder ob weitere Konfigurationen ein besseres Verhältnis zwischen Erfolgsrate und
Ausführungszeit ermöglichen.
Eine weitere Erweiterung wäre die Untersuchung zusätzlicher Fehlerklassen.
Dadurch könnte überprüft werden, ob sich das beobachtete Verhalten auch auf andere Fehler übertragen lässt.
Schließlich könnte anstelle fester Delay-Werte ein adaptiver Backoff untersucht werden, bei dem die Wartezeit
abhängig von der Anzahl bisheriger Fehlversuche angepasst wird.
Dadurch ließe sich untersuchen, ob eine dynamische Anpassung der Retry-Parameter den Kompromiss zwischen Erfolgsrate
und Ausführungszeit weiter verbessern kann.